% ============================================================
% 3_Methodology.tex
% ============================================================

\section{Methodology}
\label{sec:methodology}

We propose a distribution-driven framework for child SER consisting of four
stages: (1)~a frozen self-supervised backbone, (2)~learnable layer fusion,
(3)~temporal importance pooling (three variants compared), and (4)~an SE-MLP
classifier. The complete architecture is illustrated in Figure~\ref{fig:arch}.

\subsection{Frozen SSL Backbone}
\label{ssec:backbone}

We use WavLM Base (\texttt{wavlm-base-sv})~\cite{chen2022wavlm} as the frozen
feature extractor. All input waveforms are standardized to 16kHz mono, peak-normalized,
and truncated or padded to 4 seconds (200 frames at the 50Hz WavLM frame rate).
The backbone outputs a sequence $\mathbf{H} \in \mathbb{R}^{T \times 768}$ where
$T$ is the number of frames. The backbone's 94M parameters are frozen throughout
all experiments; we compare frozen vs.\ fully fine-tuned WavLM in dedicated
experiments (E2 series, \S\ref{ssec:wavlm_unfreeze}).

\subsection{Learnable Layer Fusion}
\label{ssec:layer_fusion}

Rather than extracting features from a single transformer layer, we compute a
learned weighted sum over all 12 WavLM layers~\cite{rozen2024layer}. Let
$\mathbf{H}^{(\ell)} \in \mathbb{R}^{T \times 768}$ be the hidden states from
layer $\ell$. The fused representation is:

\begin{equation}
\mathbf{H}^{\text{fused}} = \sum_{\ell=1}^{12} w_\ell \cdot \mathbf{H}^{(\ell)},
\quad w_\ell = \frac{\exp(\alpha_\ell)}{\sum_{k=1}^{12} \exp(\alpha_k)}
\end{equation}

where $\{\alpha_\ell\}_{\ell=1}^{12}$ are 12 learnable scalar parameters
(initialized to zero, yielding uniform weights at initialization). We compare
this against two baselines: using only the last layer (L12), and using the
single best-performing layer found via grid search on the validation set
(E5 series, \S\ref{ssec:layer_ablation}).

\subsection{Temporal Importance Pooling}
\label{ssec:pooling}

We compare three pooling strategies under strictly matched parameter budgets:

\paragraph{Mean Pooling (Baseline).}
$\mathbf{u} = \frac{1}{T}\sum_{t=1}^{T} \mathbf{h}_t$. No learnable parameters.

\paragraph{Self-Attention Pooling.}
Following~\cite{lin2017structured}, a two-layer MLP computes scalar attention
scores from $\mathbf{H}^{\text{fused}}$:

\begin{equation}
a_t = \mathbf{v}^\top \tanh(\mathbf{W}_1 \mathbf{h}_t + \mathbf{b}_1), \quad
\alpha_t = \frac{\exp(a_t)}{\sum_{j=1}^{T} \exp(a_j)}, \quad
\mathbf{u} = \sum_{t=1}^{T} \alpha_t \mathbf{h}_t
\end{equation}

where $\mathbf{W}_1 \in \mathbb{R}^{128 \times 768}$, $\mathbf{b}_1 \in \mathbb{R}^{128}$,
$\mathbf{v} \in \mathbb{R}^{128}$ (111,105 trainable parameters).

\paragraph{Prosody-Guided Pooling.}
We extract frame-aligned F0 (log-scale) and RMS energy via
\texttt{librosa}~\cite{mcfee2015librosa}, linearly interpolate to $T$ frames,
and concatenate with SSL features before computing attention:

\begin{equation}
\mathbf{p}_t = [\log\text{F0}_t, \log E_t], \quad
\tilde{\mathbf{h}}_t = [\mathbf{h}_t; \mathbf{p}_t]
\end{equation}
\begin{equation}
a_t = \mathbf{v}^\top \tanh(\mathbf{W}_1 \tilde{\mathbf{h}}_t + \mathbf{b}_1), \quad
\alpha_t = \frac{\exp(a_t)}{\sum_{j=1}^{T} \exp(a_j)}, \quad
\mathbf{u} = \sum_{t=1}^{T} \alpha_t \mathbf{h}_t
\end{equation}

Crucially, the prosody features are used only to \textit{compute} attention
weights; the pooled output is still $\sum \alpha_t \mathbf{h}_t$, keeping the
utterance representation in the SSL space. The parameter count remains
111,105---identical to self-attention pooling---because the input projection
$\mathbf{W}_1 \in \mathbb{R}^{128 \times (768+2)}$ adds only 256 extra
parameters while reducing the flattened dimension before the final projection.
We enforce exact parameter parity through a small dimension adjustment in the
linear projection, ensuring that any performance differences stem from the
prosody prior, not from additional capacity.

\subsection{SE-MLP Classifier}
\label{ssec:classifier}

The utterance vector $\mathbf{u} \in \mathbb{R}^{768}$ passes through a
Squeeze-and-Excitation MLP (SEMLP, $\sim$593K parameters): a two-layer MLP
with a channel-wise recalibration branch. The final linear layer maps to $K$
classes where $K=6$ for C-BESD and $K=4$ for FAU and IEMOCAP.

\subsection{Training Configuration}
\label{ssec:training}

All experiments use AdamW optimizer ($\text{lr}=3\times 10^{-4}$,
$\beta=(0.9,0.999)$) with cosine annealing ($T_{\text{max}}=100$), batch size
16, and early stopping with patience 15 monitored on validation WA. Speaker
splits are computed deterministically via MD5 hash with seed 42, enforcing
70/15/15 train/val/test ratios with zero-leakage assertions at runtime.

Two regularization profiles are used: \textit{default} (weight decay
$10^{-3}$, label smoothing 0.1, no pooling dropout) for C-BESD and IEMOCAP;
\textit{fau} (weight decay $5\times 10^{-3}$, label smoothing 0.15, pooling
dropout 0.3, gradient clipping 1.0) for FAU Aibo to combat overfitting on the
smaller labeled subset. Unless otherwise stated, every experiment is repeated
with three fixed random seeds (42, 123, 456) and we report mean $\pm$ standard
deviation.

\subsection{Dataset-Specific Processing}
\label{ssec:datasets}

\paragraph{C-BESD (6-class).} We collect all 4,179 utterances from 6 emotion
categories (ANGER, DISGUST, FEAR, HAPPY, NEUTRAL, SAD) across 70 unique
children. Child identity is extracted from the filename pattern
(\texttt{\{child\_id\}.EF/TF/EM/TM\_\{session\}\_\{emotion\}\_\{index\}.wav})
using regex matching on the leading integer: `1.EF\_12 Angry\_1.wav'
$\rightarrow$ child `C01'. This extraction is critical: prior work that used
the full session prefix as speaker ID risked placing different sessions of the
same child in different splits, creating undetected data leakage. Our
extraction yields 70 child identifiers (vs.\ 237 session-level strings),
ensuring true speaker-disjoint partitioning.

\paragraph{FAU Aibo (4-class).} The corpus contains 18,216 spontaneous
utterances from 51 German children (aged 10--13, two schools: Mont, Ohm)
recorded during child--robot interaction. We use the chunk-level 5-class
annotation (A=Anger, E=Emphatic, N=Neutral, P=Positive, R=Rest) and remap to
a 4-class taxonomy by merging Anger and Emphatic into \textit{Angry}
(A+E$\rightarrow$angry, 5,093 utterances), treating Positive as \textit{Happy}
(889), Neutral as \textit{Neutral} (10,967), and Rest as \textit{Sad} (1,267).
This remapping aligns FAU's label space with the 4-class subsets of C-BESD and
IEMOCAP, enabling cross-corpus comparison. Speaker identity is extracted from
the filename as \texttt{\{school\}\_\{child\_id\}} (51 unique speakers, verified
against the official age.txt metadata).

\paragraph{IEMOCAP (4-class).} We retain the four most frequent emotion
categories (angry, happy, neutral, sad) following standard practice, discarding
1,269 utterances with labels outside this set (fear, disgust, surprise,
excited, frustrated). Speaker identity is extracted from the first 6 characters
of the filename identifying the session and gender (e.g., `Ses01F'), yielding
10 unique speakers.

\paragraph{Preprocessing Pipeline.} All audio is resampled to 16kHz mono,
peak-normalized, and truncated/padded to 4 seconds. Speaker splitting uses
MD5-hash deterministic assignment at 70/15/15 ratios with seed 42, verified by
runtime zero-overlap assertions. No speaker appears in more than one split.
\end{document}
